\section*{الفصل \ref{ch:g11:enzymes-and-phenotype} --- الإنزيمات والنمط الظاهري}

\begin{solution}{exo:g11:enzymes-and-phenotype:1}
\omterm{def:g11:enzymes-and-phenotype:phenotype}{النمط الوراثي}: مجموعة \omterm{def:g10:universal-dna:gene}{الأليلات} المحمولة. \omterm{def:g11:enzymes-and-phenotype:phenotype}{والنمط الظاهري}: مجموعة
الخصائص القابلة للمشاهدة، على السلالم الجزيئي والخلوي والعضوي.
\end{solution}

\begin{solution}{exo:g11:enzymes-and-phenotype:2}
جيب \omterm{def:g11:enzymes-and-phenotype:enzyme}{الإنزيم} الذي ترتبط فيه \omterm{def:g11:enzymes-and-phenotype:enzyme}{الركيزة} وتتحول. أما النوعية: فهي أن الموقع
يوافق \omterm{def:g11:enzymes-and-phenotype:enzyme}{ركيزة} واحدة (أو عائلة واحدة من الركائز) ويحفّز تفاعلًا واحدًا.
\end{solution}

\begin{solution}{exo:g11:enzymes-and-phenotype:3}
قرب \qty{37}{\celsius}؛ وعند \qty{45}{\celsius} يبقى نحو الثلث.
\end{solution}

\begin{solution}{exo:g11:enzymes-and-phenotype:4}
جزيئيًا: التيروزيناز غائب أو خامل، ولا يصنع ميلانين. وخلويًا: \omterm{def:g10:cells-common-unit:cell}{خلايا}
الصباغ موجودة لكنها خالية من الميلانين. وعضويًا: شعر أبيض، وجلد
وعينان شديدا الشحوب، وحساسية للضوء ولحروق الشمس.
\end{solution}

\begin{solution}{exo:g11:enzymes-and-phenotype:5}
طيّة البيبسين لا تعمل إلا في حمض قوي (pH 2)؛ والمعي متعادل إلى قلوي
قليلًا، حيث يبين المنحني نشاطًا صفريًا.
\end{solution}

\begin{solution}{exo:g11:enzymes-and-phenotype:6}
\omterm{def:g10:universal-dna:gene}{الأليل} العامل الواحد ينتج نصف \omterm{def:g11:enzymes-and-phenotype:enzyme}{الإنزيم} الطبيعي، لكن جزيء \omterm{def:g11:enzymes-and-phenotype:enzyme}{الإنزيم} يعمل
آلاف المرات في الثانية: فنصف الكمية لا يزال يحول ما يكفي من التيروزين
لصنع صباغ طبيعي. \omterm{def:g11:enzymes-and-phenotype:phenotype}{فالنمط الظاهري} مشبع؛ ولا يظهر إلا مع \omterm{def:g10:universal-dna:gene}{أليلين} مخفقين.
\end{solution}

\begin{solution}{exo:g11:enzymes-and-phenotype:7}
مهقهما آتٍ من حصر في مورّثتين مختلفتين (إنزيمان مختلفان من السبيل، أو
من نقل الميلانين). فيرث الطفل أليلًا عاملًا لكل \omterm{def:g10:universal-dna:gene}{مورّثة} من الوالد الذي
حصره في مكان آخر، فيعمل الإنزيمان ويصنع الصباغ.
\end{solution}

\begin{solution}{exo:g11:enzymes-and-phenotype:8}
يتراكم الفينيل ألانين؛ ويفتقر التيروزين (وما يشتق منه). ولما كان
الميلانين يصنع من التيروزين، فنقصه يقلل الصباغ: فالشعر والجلد الفاتحان
شائعان في بيلة الفينيل كيتون غير المعالجة.
\end{solution}

\begin{solution}{exo:g11:enzymes-and-phenotype:9}
حيث يكون الحليب غذاءً أساسيًا، يوفر هضمه الطاقة والكالسيوم بلا مرض،
فيكون \omterm{def:g10:universal-dna:gene}{الأليل} الباقي ميزة. وحيث لا يشرب البالغون حليبًا لا يغير \omterm{def:g10:universal-dna:gene}{الأليل}
شيئًا؛ فهو «أصح» في المحيط الذي يوفر \omterm{def:g11:enzymes-and-phenotype:enzyme}{الركيزة} وحده.
\end{solution}

\begin{solution}{exo:g11:enzymes-and-phenotype:10}
عند التيروزيناز: جزيئيًا، لا تلائم \omterm{def:g11:enzymes-and-phenotype:enzyme}{الركيزة} بعد ذلك، ولا ميلانين؛
وخلويًا، \omterm{def:g10:cells-common-unit:cell}{خلايا} صباغ فارغة؛ وعضويًا، مهق. وعند اللاكتاز: لاكتوز غير
مهضوم؛ \omterm{def:g10:cells-common-unit:cell}{وخلايا} معوية محرومة من السكر والقولون يخمّره؛ وعدم احتمال.
\end{solution}

\begin{solution}{exo:g11:enzymes-and-phenotype:11}
أطرافها تبقى أدفأ من \qty{35}{\celsius}، فلا يكاد \omterm{def:g11:enzymes-and-phenotype:enzyme}{الإنزيم} يعمل، وتكون
الهريرة فاتحة كلها. أما في الخارج شتاءً فتبرد الأذنان والكفوف، ويصير
\omterm{def:g11:enzymes-and-phenotype:enzyme}{الإنزيم} فعالًا، ويأتي الفرو النامي هناك داكنًا.
\end{solution}

\begin{solution}{exo:g11:enzymes-and-phenotype:12}
جزيئيًا: نصف الهيموغلوبين من الصورة المنجلية. وخلويًا: عند أكسجين
طبيعي تبقى \omterm{def:g10:cells-common-unit:cell}{الخلايا} مدورة؛ وعند أكسجين منخفض جدًا يتمنجل بعضها.
وعضويًا: سليم في الحياة العادية، ومعرض للخطر على الارتفاعات أو عند
الغوص --- ومحمي من الملاريا. فالمحيط (الأكسجين والملاريا) يقرر هل نمط
الحامل الظاهري عبء أم ميزة.
\end{solution}

\begin{solution}{exo:g11:enzymes-and-phenotype:13}
عند \qty{41}{\celsius} تحفظ \omterm{def:g11:enzymes-and-phenotype:enzyme}{الإنزيمات} معظم نشاطها (فالمنحني قرب ذروته)؛
وعند \qty{43}{\celsius} يهبط النشاط هبوطًا حادًا إذ تبدأ \omterm{def:g10:chemistry-of-life:families}{البروتينات}
بالانتشار، وإذا مضى الانتشار كان غير قابل للرد: \omterm{def:g11:enzymes-and-phenotype:enzyme}{فالإنزيمات} تهدم لا
تبطأ.
\end{solution}

\begin{solution}{exo:g11:enzymes-and-phenotype:14}
متقدرات وشعيرات دموية أكثر لكل ليف عضلي؛ وقلب أكبر بحجم نفضي أكبر؛
وهيموغلوبين أكثر في لتر الدم؛ ومخازن \omterm{def:g10:exercise-and-energy:glycogen}{غليكوجين} أكبر. وكل منها \omterm{def:g11:enzymes-and-phenotype:phenotype}{نمط
ظاهري} ينتجه \omterm{def:g11:enzymes-and-phenotype:phenotype}{النمط الوراثي} نفسه في محيط مختلف: فالصفات «الوراثية» هي
المدى الذي يتيحه \omterm{def:g11:enzymes-and-phenotype:phenotype}{النمط الوراثي}، لا قيمة ثابتة.
\end{solution}

\begin{solution}{exo:g11:enzymes-and-phenotype:15}
\omterm{def:g11:enzymes-and-phenotype:phenotype}{النمط الوراثي} يثبّت أي \omterm{def:g11:enzymes-and-phenotype:enzyme}{إنزيم} يصنع وكيف يسلك: فلا \omterm{def:g11:enzymes-and-phenotype:enzyme}{إنزيم} 1 (بيلة الفينيل
كيتون)، ولا تيروزيناز (المهق)، أو تيروزيناز حساس للحرارة (السيامي).
والمحيط لا يستطيع تغيير النتيجة إلا في حدود ما يتيحه \omterm{def:g11:enzymes-and-phenotype:enzyme}{الإنزيم}: فالغذاء
يزيل \omterm{def:g11:enzymes-and-phenotype:enzyme}{الركيزة} وينقذ الدماغ، لكنه لا يستطيع صنع \omterm{def:g11:enzymes-and-phenotype:enzyme}{الإنزيم} المفقود؛ والبرد
يجعل \omterm{def:g11:enzymes-and-phenotype:enzyme}{الإنزيم} السيامي يعمل، لكن ما من حرارة تجعل \omterm{def:g11:enzymes-and-phenotype:enzyme}{إنزيم} الأمهق يعمل؛
وما من شيء في المحيط يعطي الأمهق صباغه. \omterm{def:g11:enzymes-and-phenotype:phenotype}{فالنمط الوراثي} يضع المدى؛
والمحيط يختار داخله.
\end{solution}

\begin{solution}{pb:g11:enzymes-and-phenotype:1}
\textbf{1.} البدن \qty{38}{\celsius}: نحو 2\%، فاتح. والأرجل
\qty{36}{\celsius}: نحو 20\%، فاتح. والكفوف \qty{33}{\celsius}: 80\%،
داكن. والأذنان والخطم \qty{31}{\celsius}: أكثر من 90\%، داكن.

\textbf{2.} بدن وأرجل فاتحة، وكفوف وأذنان وخطم (وذنب) داكنة: أي النمط
السيامي.

\textbf{3.} 100\% في كل مكان: أي قط متصبغ تصبغًا منتظمًا.

\textbf{4.} في الرحم يكون كل جزء من الهريرة عند \qty{38}{\celsius}،
\omterm{def:g11:enzymes-and-phenotype:enzyme}{والإنزيم} خامل في كل مكان، فلا يرسَّب صباغ قبل الولادة؛ وتدكن النقاط
حين تبرد الأطراف بعد الولادة.

\textbf{5.} قرب \qty{35.5}{\celsius}، حيث يعبر النشاط نسبة 30\%:
فالجزء النازل الحاد من المنحني بين 33 و\qty{37}{\celsius} يثبّت الحد.

\textbf{6.} داكن: فعند \qty{30}{\celsius} يكون \omterm{def:g11:enzymes-and-phenotype:enzyme}{الإنزيم} فعالًا تمامًا
في الشعرات النامية.

\textbf{7.} فاتح: فعند \qty{38}{\celsius} يكون \omterm{def:g11:enzymes-and-phenotype:enzyme}{الإنزيم} خاملًا.

\textbf{8.} يرسَّب الصباغ في الشعرة وهي تنمو ولا يمكن نزعه ولا إضافته
بعد ذلك؛ ولا يتغير اللون إلا حين تتساقط الشعرات وتستبدل بشعرات جديدة
نمت عند الحرارة الموضعية.

\textbf{9.} \omterm{def:g11:enzymes-and-phenotype:phenotype}{النمط الوراثي} \omterm{def:g10:cells-common-unit:cell}{لخلايا} الرقعة هو هو قبل وبعد، وهو نفسه في
سائر البدن؛ ولم يتغير إلا حرارة \omterm{def:g11:enzymes-and-phenotype:enzyme}{الإنزيم}، ومعها الصباغ. فالفعل واقع على
نشاط \omterm{def:g11:enzymes-and-phenotype:enzyme}{الإنزيم}، وهو حدث جزيئي.

\textbf{10.} \omterm{def:g10:cells-common-unit:cell}{خلايا} القط كلها تنحدر من بيضة واحدة وتحمل \omterm{def:g10:universal-dna:gene}{الأليلات} نفسها؛
ثم إن تجارب الرقعة المحلوقة تغير لون منطقة من دون تغيير خلاياها ---
\omterm{def:g10:universal-dna:gene}{فالأليل} لا يبدله كيس ثلج.

\textbf{11.} جزيئيًا: تيروزيناز لا يعمل إلا دون نحو \qty{35}{\celsius}.
وخلويًا: \omterm{def:g10:cells-common-unit:cell}{خلايا} صباغ مملوءة ميلانينًا في الأطراف، فارغة في البدن.
وعضويًا: قط فاتح ذو نقاط داكنة.

\textbf{12.} \omterm{def:g10:chemistry-of-life:families}{الحمض الأميني} يقع حيث يثبّت الطيّة؛ والبديل يمسك الطيّة
عند الحرارة المنخفضة لكنه يدعها ترتخي عند درجات قليلة أعلى، فيفقد
\omterm{def:g11:enzymes-and-phenotype:enzyme}{الموقع الفعال} شكله قبل أن يفقده موقع \omterm{def:g11:enzymes-and-phenotype:enzyme}{الإنزيم} العادي.

\textbf{13.} \omterm{def:g11:enzymes-and-phenotype:enzyme}{إنزيم} \omterm{def:g10:universal-dna:gene}{الأليل} العادي فعال تمامًا عند حرارة الجسم؛ ونصف
الكمية الطبيعية أكثر من كاف لصنع الصباغ في كل مكان. \omterm{def:g10:universal-dna:gene}{فالأليل} السيامي
متنح أمامه.

\textbf{14.} جزيئيًا: لا \omterm{def:g11:enzymes-and-phenotype:enzyme}{إنزيم} عامل عند أي حرارة، مقابل \omterm{def:g11:enzymes-and-phenotype:enzyme}{إنزيم} يعمل
شرطيًا. وخلويًا: لا ميلانين في أي مكان، مقابل ميلانين في المناطق
الباردة. وعضويًا: قط أبيض كله بعينين ورديتين، مقابل النمط ذي النقاط.
والبرد لا يغير شيئًا عند الأمهق.

\textbf{15.} أنها تحمل أليلًا حساسًا للحرارة \omterm{def:g10:universal-dna:gene}{للمورّثة} نفسها؛ وعلى
الأرجح استبدالًا مماثلًا يزعزع \omterm{def:g11:enzymes-and-phenotype:enzyme}{الإنزيم}، نشأ مستقلًا في كل \omterm{def:g10:biodiversity-scales:species}{نوع}.

\textbf{16.} عشرين ضعفًا؛ \omterm{def:g11:enzymes-and-phenotype:enzyme}{والإنزيم} 1، وهو الذي يحول الفينيل ألانين إلى
تيروزين.

\textbf{17.} فائض الفينيل ألانين يتداخل مع \omterm{def:g10:cells-common-unit:cell}{خلايا} الدماغ النامي بعينها
(في إمدادها \omterm{def:g10:chemistry-of-life:families}{بالحموض الأمينية} الأخرى وفي نضجها)؛ أما \omterm{def:g10:cells-common-unit:cell}{خلايا} الجلد
فتحتمله. فالعيب الجزيئي واحد في كل مكان؛ أما العاقبة الخلوية فلا.

\textbf{18.} خمس \num{1200} هو \qty{240}{\micro mol/L}، وهو دون
\num{360}: فآمن. ومن دون \omterm{def:g11:enzymes-and-phenotype:enzyme}{الإنزيم} يتبع مستوى الدم الوارد ببساطة، فيضبط
النمطَ الظاهري إمدادُ \omterm{def:g11:enzymes-and-phenotype:enzyme}{الركيزة} --- أي المحيط --- لا \omterm{def:g11:enzymes-and-phenotype:enzyme}{الإنزيم} المفقود
وحده.

\textbf{19.} الفينيل ألانين لازم لبناء بروتينات الجسم ولا يستطيع الجسم
تركيبه؛ فإذا انعدم توقف النمو. فالوارد يجب أن يكفي تركيب \omterm{def:g11:gene-expression:protein}{البروتين} ولا
يزيد، وهو ما يقتضي قياسًا وضبطًا.

\textbf{20.} \omterm{def:g11:enzymes-and-phenotype:phenotype}{النمط الوراثي} ثبّت \omterm{def:g11:enzymes-and-phenotype:enzyme}{الإنزيم} (حساسًا للحرارة عند القط،
غائبًا عند الطفل)؛ والمحيط ضبط شروطه (حرارة الجلد؛ وفينيل ألانين
الغذاء)؛ \omterm{def:g11:enzymes-and-phenotype:phenotype}{والنمط الظاهري} حصيلة \omterm{def:g11:enzymes-and-phenotype:phenotype}{النمط الوراثي} معبَّرًا عنه في محيط معطى.
\end{solution}
